% !TeX root = ../../Book3.tex
% 纲要标题：### 10.1 赋值与赋值环
% 纲要位置：工作记录/计划纲要.md，第 2230 行。
% 纲要细目（写作范围）：
% * 赋值群及赋值环的整除性质
% * 一般赋值环与离散赋值环的区别
% * 分式域和局部性
% * 赋值的秩与赋值环的 Krull 维数：用凸子群与素理想的反包含对应说明维数，区分秩与有理秩，并使用第9章已建立的素链语言

\section{赋值与赋值环}
\label{b3:sec:1001}

整数中可以比较两个数含有多少个素因子 \(p\)，有理函数中也可以比较它们在一点的零点或极点阶数。这些阶数在乘法中相加，例如 \(p^r p^s=p^{r+s}\)；相加时则可能因最低阶项抵消而升高阶数。赋值把这种计算写成有序群中的加法。允许一般的有序群，才能看清哪些性质来自整除的全序性，哪些性质另外需要离散性或 Noether 性。


\subsection{赋值群与赋值环的整除性质}
\label{b3:subsec:1001:01}

\begin{definition}\label{b3:def:valuation}
设 \(\Gamma\) 为带有全序的加法交换群。若对所有 \(\alpha,\beta,\gamma\in\Gamma\) 有
\(\alpha\le\beta\Rightarrow\alpha+\gamma\le\beta+\gamma\)，则称 \(\Gamma\) 为\term[quanxu-jiaohuanqun]{全序交换群}[totally ordered abelian group]。
设 \(K\) 为域，\(\Gamma\) 为全序交换群。给定满射
\(v:K^*\rightarrow\Gamma\)，并规定 \(v(0)=+\infty\)。若对非零 \(a,b\in K\) 满足第一式，且对任意 \(a,b\in K\) 满足第二式：
\[
v(ab)=v(a)+v(b),\qquad
v(a+b)\ge\min\bigl\{v(a),v(b)\bigr\},
\]
则称 \(v\) 为 \(K\) 上取值于 \(\Gamma\) 的\term[fuzhi]{赋值}[valuation]。其中 \(+\infty\) 大于每个群元素，且
\(\gamma+(+\infty)=+\infty\)。称 \(\Gamma\) 为赋值群。
\end{definition}
\symidx{\(v:K^*\to\Gamma\)：域的赋值及其值群}

本书约定 \(\Gamma\) 就是实际取到的值群，所以要求 \(v\) 满射。若采用允许映射不满射的定义，只须把所给有序群换成子群 \(v(K^*)\)，就得到这里的约定；赋值不等式及相应赋值环都不改变。

\ProseParagraphBreak
由第一条公理有 \(v(1)=0\)、\(v(a^{-1})=-v(a)\)。全序群无挠，故由
\(2v(-1)=0\) 得 \(v(-1)=0\)。如果 \(v(a)<v(b)\)，则
\(v(a+b)=v(a)\)：否则把 \(a=(a+b)-b\) 代入三角不等式便矛盾。
这说明最低阶项只有在至少两项同阶时才可能互相抵消。例如对任意域 \(k\) 及不定元 \(t\)，\(t^2+t^5=t^2(1+t^3)\)，括号内在 \(t=0\) 处的值为一，所以和在零点处的阶仍为二。

\begin{definition}\label{b3:def:valuation-ring}
设 \(K\) 为域，\(v:K^*\rightarrow\Gamma\) 为取值于全序交换群 \(\Gamma\) 的赋值。定义子环
\[
V_v=\{0\}\cup\bigl\{a\in K^*\mid v(a)\ge0\bigr\},
\]
并称 \(V_v\) 为 \(v\) 的\term[fuzhihuan]{赋值环}[valuation ring]。
等价地，若整环 \(V\) 的分式域为 \(K\)，且对每个 \(x\in K^*\) 都有
\(x\in V\) 或 \(x^{-1}\in V\)，则称 \(V\) 为 \(K\) 的赋值环。允许 \(V=K\)，这对应零赋值群，称为平凡赋值。
\end{definition}
\symidx{\(V_v\)：赋值 \(v\) 的赋值环}

说明两种说法确实等价。赋值不等式保证 \(V_v\) 对加法、减法、乘法封闭，且含 \(1\)；全序性给出倒数条件。
反过来，若 \(V\) 满足倒数条件，乘一个 \(V\) 的单位不会改变整除关系。于是把只差一个单位的非零元素看成具有同一个阶，即把满足 \(a/b\in V^*\) 的 \(a,b\) 识别为同一类。所得商群是 \(\Gamma=K^*/V^*\)，它忘掉单位因子，保留整除的比较；把商群运算写成加法，并规定
\[
[a]\le[b]\quad\Longleftrightarrow\quad b/a\in V.
\]
换乘单位不改变此条件；若两个方向同时成立，商为单位，故两类相等。
倒数条件给出全序，乘法给出与平移的相容性。设 \(v(a)=[a]\)，若
\(v(a)\le v(b)\)，则 \((a+b)/a=1+b/a\in V\)，于是赋值不等式成立。

\begin{proposition}\label{b3:prop:valuation-ideal-comparison}
设 \(v:K^*\rightarrow\Gamma\) 为域 \(K\) 上的赋值，\(V_v\) 为其赋值环。则 \(V_v\) 的任意两个理想可以按包含关系比较。因此每个有限生成理想都是主理想。
对非零元素 \(a,b\in V_v\)，有
\(a\mid b\Longleftrightarrow v(a)\le v(b)\)。
\end{proposition}
\begin{proof}
末一等价式就是 \(b/a\in V_v\)。全序性因此先保证两个非零元素生成的主理想可以比较，再由此比较任意理想。若理想 \(I\) 不包含于 \(J\)，取
\(a\in I\setminus J\)。任取 \(0\ne b\in J\)，不能有 \(a\in bV_v\)，故必有
\(b\in aV_v\subseteq I\)，于是 \(J\subseteq I\)。

若 \(I=(a_1,\ldots,a_r)\ne0\)，在非零生成元中取赋值最小的 \(a_j\)。每个非零 \(a_i\) 都满足 \(v(a_j)\le v(a_i)\)，故 \(a_j\mid a_i\)，从而 \(I=a_jV_v\)。零理想本身由零生成，所以每个有限生成理想都是主理想。
\end{proof}

\subsection{一般赋值环与离散赋值环}
\label{b3:subsec:1001:02}

\begin{definition}\label{b3:def:discrete-valuation-ring}
设 \(K\) 为域，\(v:K^*\rightarrow\Gamma\) 为域 \(K\) 上的赋值。若赋值群 \(\Gamma\) 与通常排序的 \(\mathbb Z\) 有序同构，则称相应赋值环 \(V_v\) 为
离散赋值环\termidx[lisan-fuzhi-huan]{离散赋值环}，缩写为 DVR。本书的 DVR 不是域。
\end{definition}

“离散”说的是赋值群由整数阶组成，并不是说极大理想的幂所定义的拓扑为离散拓扑。

\ProseParagraphBreak
对素数 \(p\)，将非零有理数写成 \(p^n a/b\)，其中 \(a,b\) 为整数且 \(p\nmid ab\)。指数 \(n\) 唯一，定义
\(v_p(p^na/b)=n\)，所得赋值环是 \(\mathbb Z_{(p)}\)。

\ProseParagraphBreak
设 \(k\) 为域，\(t\) 为不定元。每个非零 \(f\in k[t]\) 唯一写成 \(f=t^ru\)，其中 \(r\ge0\)、\(u\in k[t]\) 且 \(u(0)\ne0\)。记 \(\nu_t(f)=r\)，即 \(t\) 在 \(f\) 中作为因子的重数。对 \(f,g\in k[t]\setminus\{0\}\)，定义
\[
\operatorname{ord}_0(f/g)=\nu_t(f)-\nu_t(g).
\]
因子重数与多项式次数不同，例如 \(\operatorname{ord}_0(t+1)=0\)。这个赋值的赋值环为 \(k[t]_{(t)}\)。两例中乘法阶数相加、加法阶数不下降，都可直接由提取相应素因子的幂验证。

\ProseParagraphBreak
一般赋值未必离散。在 \(\mathbb Q\) 中，任何正值 \(q\) 下方仍有正值 \(q/2\)，所以正值没有最小元。若一个赋值以 \(\mathbb Q\) 为实际值群，正赋值元素连同零组成的理想就不能有限生成：有限生成会给出一个赋值最小的生成元，其正值记为 \(q\)；满射性却给出赋值为 \(q/2\) 的元素，它属于这个理想而不能被该生成元整除。

\ProseParagraphBreak
这样的赋值可如下构造。设 \(k\) 为域，\(t\) 为不定元，在一个固定扩域中相容选取 \(t^{1/n!}\)，并令
\(K=\bigcup_{n\ge1}k(t^{1/n!})\)；
取 \(v(t^{1/n!})=1/n!\)，在 \(k(t^{1/n!})\) 中把以 \(t^{1/n!}\) 为变量的零点阶除以 \(n!\)。这些赋值在嵌套子域上相容，因而定义了 \(K\) 上以 \(\mathbb Q\) 为值群的赋值。
每个有理数的分母都整除某个 \(n!\)，所以各子域的值群 \(\dfrac1{n!}\mathbb Z\) 的并恰为 \(\mathbb Q\)。
因此所得赋值环不是 Noether 环。

\subsection{分式域与局部性}
\label{b3:subsec:1001:03}

\begin{proposition}\label{b3:prop:valuation-local-normal}
设 \(v:K^*\rightarrow\Gamma\) 为域 \(K\) 上的赋值，\(V_v\) 为其赋值环。则 \(V_v\) 的分式域是 \(K\)，它是局部整环，唯一极大理想为
\[
\mathfrak m_v=\{0\}\cup\bigl\{a\in K^*\mid v(a)>0\bigr\},
\]
单位恰为赋值零的元素。每个赋值环都在 \(K\) 中整闭。
\end{proposition}
\symidx{\(\mathfrak m_v\)：赋值环中正赋值元素组成的极大理想}
\begin{proof}
任意 \(x\in K^*\) 或属于 \(V_v\)，或为 \(V_v\) 中非零元素的倒数，故分式域为 \(K\)。
赋值公理说明 \(\mathfrak m_v\) 是理想；非零 \(a\in V_v\) 可逆恰当 \(v(a)=0\)，所以非单位全体恰是该真理想。这证明局部性。

整闭性要求排除负赋值元素满足 \(V_v\) 系数首一方程的可能。若 \(x\in K\) 满足首一方程 \(x^n+a_1x^{n-1}+\cdots+a_n=0\)、\(a_i\in V_v\)，而 \(v(x)<0\)，则 \(x^n\) 的赋值严格低于其余各个非零项的赋值：系数的赋值非负，较低次的 \(x\) 的幂反而具有较高赋值。这唯一的最低阶项不能被抵消。把这一点写成理想中的等式，除以 \(x^n\) 得
\(1+a_1x^{-1}+\cdots+a_nx^{-n}=0\)。除首项外每项均属于 \(\mathfrak m_v\)，于是 \(1\in\mathfrak m_v\)，矛盾。因此 \(x\in V_v\)。
\end{proof}

局部性不保证是赋值环。例如对域 \(k\)，在 \(R=k[x,y]_{(x,y)}\) 中，主理想 \((x)\) 与 \((y)\) 不能比较：模掉 \((y)\) 后得到 \(k[x]_{(x)}\)，其中 \(x\) 非零，故 \(x\notin(y)\)；交换两变量，同样有 \(y\notin(x)\)。所以 \(R\) 不是赋值环；这两个不包含关系也说明 \(x/y\) 与 \(y/x\) 均不在 \(R\) 中。多项式环 \(k[x,y]\) 整闭，由\cref{b3:cor:integrally-closed-localization}，\(R\) 也整闭，因此整闭局部环仍不必是赋值环。

\subsection{赋值的秩与赋值环的 Krull 维数}
\label{b3:subsec:1001:04}

设 \(\Gamma\) 为全序交换群，\(H\) 为其子群。若对所有 \(\gamma\in\Gamma\) 和 \(h\in H\)，
\(0\le\gamma\le h\) 均蕴含 \(\gamma\in H\)，则称 \(H\) 为\term[tuzhiqun]{凸子群}[convex subgroup]。
凸子群按包含关系组成全序：若 \(H_1\not\subseteq H_2\)，取正的
\(h\in H_1\setminus H_2\)，则 \(H_2\) 的每个正元素均小于 \(h\)，故属于 \(H_1\)。
赋值的\term[fuzhi-zhixu]{秩}[rank of a valuation]定义为其赋值群中有限严格凸子群链的长度上确界。它测量凸子群能分出多少层，所数的不是群生成元的个数。赋值的\term[youli-zhixu]{有理秩}[rational rank]则定义为
\[
\operatorname{rat.rank}(v)
=\dim_{\mathbb Q}(\Gamma\otimes_{\mathbb Z}\mathbb Q),
\]
它测量值群有理化后的向量空间维数。
\symidx{\(\operatorname{rat.rank}(v)\)：赋值 \(v\) 的有理秩}

\begin{samepage}
以按字典序排序的 \(\Gamma=\mathbb Z^2\) 为例，先比较第一坐标，第一坐标相同时才比较第二坐标。因此 \((0,n)\) 对任意整数 \(n\) 都小于 \((1,0)\)。子群 \(H=\{0\}\times\mathbb Z\) 是凸的：夹在零与 \(H\) 的正元素之间的元素，第一坐标只能为零。非零凸子群若只含第一坐标为零的元素，凸性便迫使它包含 \((0,1)\)，故等于 \(H\)；若含第一坐标为正的元素，其正整数倍最终超过任意给定正元素，凸性便使该子群等于全群。

\ProseParagraphBreak
对于值群为 \(\mathbb Z^2\) 的赋值，令 \(\mathfrak p\) 由零元素以及赋值第一坐标为正的元素组成。它只检测第一层阶数，忽略第二坐标。凸子群与素理想之间的对应排列为
\[
\begin{array}{ccccc}
0&\subsetneq&\{0\}\times\mathbb Z&\subsetneq&\mathbb Z^2\\
\updownarrow&&\updownarrow&&\updownarrow\\
\mathfrak m_v&\supsetneq&\mathfrak p&\supsetneq&(0).
\end{array}
\]
保留的凸子群越大，对应素理想越小。一般值群也有这样的反向对应。
\end{samepage}

\begin{proposition}\label{b3:prop:valuation-rank-dimension}
设 \(v:K^*\rightarrow\Gamma\) 为域 \(K\) 上的赋值，\(V_v\) 为其赋值环。则 \(\Gamma\) 的凸子群与 \(V_v\) 的素理想有反包含对应
\[
H\longmapsto\mathfrak p_H
=\{0\}\cup\bigl\{a\in V_v\setminus\{0\}\mid v(a)>h\;\text{对所有}\;h\in H\bigr\}.
\]
特别地，赋值的秩等于赋值环的 Krull 维数。
\end{proposition}
\begin{proof}
对非负群元素，属于 \(H\) 或大于 \(H\) 中所有元素，二者必居其一；这是凸性的直接结果。
所以 \(V_v\setminus\mathfrak p_H\) 恰由赋值属于 \(H\) 的非零元素组成。
这个集合对乘法封闭，而 \(\mathfrak p_H\) 对相加、乘以 \(V_v\) 元素封闭，故为素理想。

反过来，给定素理想 \(\mathfrak p\)，令
\(S=\bigl\{v(a):a\in V_v\setminus\mathfrak p\bigr\}\subseteq\Gamma_{\ge0}\)。
它对加法封闭。若 \(0\le\gamma\le s=v(a)\in S\)，取 \(b\in V_v\) 使 \(v(b)=\gamma\)；
因 \(b\mid a\)，\(b\in\mathfrak p\) 会迫使 \(a\in\mathfrak p\)，所以 \(\gamma\in S\)。
于是 \(S\) 向下封闭；若 \(s,t\in S\) 且 \(s\ge t\)，则 \(0\le s-t\le s\)，故 \(s-t\in S\)。所以 \(H=S-S=S\cup(-S)\) 是凸子群，其非负部分恰为 \(S\)。
因此 \(\mathfrak p=\mathfrak p_H\)，两构造互逆。包含的反转由公式直接看出。
有限严格链由此互相对应，长度不变。
\end{proof}

实数加法群的任意非零有序子群都只有两个凸子群：若凸子群含正元素 \(h\)，则 \(h\) 的正整数倍最终超过任意给定正元素，凸性使它等于全群。因此通常排序的 \(\mathbb Z\)、\(\mathbb Q\) 和 \(\mathbb Z+\sqrt2\mathbb Z\) 的秩均为一。前两个群有理化后都是一维 \(\mathbb Q\) 向量空间，有理秩为一；第三个群中的 \(1,\sqrt2\) 在 \(\mathbb Q\) 上线性独立，有理秩为二。字典序 \(\mathbb Z^2\) 的凸子群链则有两层，秩为二，其有理秩也为二。

\ProseParagraphBreak
这四个值群中，只有通常排序的 \(\mathbb Z\) 对应本书的 DVR。\(\mathbb Q\) 和 \(\mathbb Z+\sqrt2\mathbb Z\) 虽然秩一，却不与 \(\mathbb Z\) 有序同构；字典序 \(\mathbb Z^2\) 虽有最小正值 \((0,1)\)，秩仍为二。因此秩一或有最小正值都不能单独替代离散赋值环的定义。上面的 \(\mathbb Q\) 值群例子是一维整闭局部环，却不是 DVR；下一节将考察一维局部整环再满足 Noether 性和整闭性时的结构。
